\begin{tikzpicture}[x=1mm,y=1mm]
\tikzset{
 fn/.style={draw=cinzamedio,fill=white,align=center,font=\footnotesize,inner xsep=1mm,inner ysep=1mm,text width=22mm,minimum height=13mm,execute at begin node=\setstretch{1}},
 fp/.style={fn,fill=cinzaclaro,minimum height=14mm},
 material/.style={-{Stealth[length=1.7mm]},draw=cinzaescuro,line width=.8pt},
 energia/.style={-{Stealth[length=1.5mm]},draw=cinzamedio,dashed,line width=.6pt},
 sinal/.style={-{Stealth[length=1.5mm]},draw=cinzamedio,dash pattern=on 3pt off 2pt on .6pt off 2pt,line width=.6pt},
 rel/.style={-{Stealth[length=1mm]},draw=cinzamedio,line width=.45pt},
 porta/.style={align=left,font=\footnotesize,inner sep=0,execute at begin node=\setstretch{1}},
 legenda/.style={font=\footnotesize,anchor=west,execute at begin node=\setstretch{1}}
}

\fill[cinzaclaro,opacity=.45] (26,8) rectangle (52,-75);
\node[fp] (p0) at (39,0) {{\scriptsize P4}\\\textbf{Acomodar louça}};
\node[fn] (e00) at (39,-19) {{\scriptsize E10}\\\textbf{Receber louça}};
\node[fn] (e01) at (39,-34) {{\scriptsize E11}\\\textbf{Ajustar apoios}};
\node[fn] (e02) at (39,-49) {{\scriptsize E12}\\\textbf{Posicionar louça}};
\draw[energia] (39,-82)--(39,-75);
\node[porta,anchor=north,align=center,text width=25mm,font=\scriptsize] at (39,-84) {Trabalho do operador};
\fill[cinzaclaro,opacity=.45] (53,8) rectangle (79,-75);
\node[fp] (p1) at (66,0) {{\scriptsize P1}\\\textbf{Imobilizar louça}};
\node[fn] (e10) at (66,-19) {{\scriptsize E01}\\\textbf{Travar apoios}};
\node[fn] (e11) at (66,-34) {{\scriptsize E02}\\\textbf{Distribuir pressão}};
\node[fn] (e12) at (66,-49) {{\scriptsize E03}\\\textbf{Restringir movimento}};
\draw[energia] (66,-82)--(66,-75);
\node[porta,anchor=north,align=center,text width=25mm,font=\scriptsize] at (66,-84) {Trabalho do operador};
\fill[cinzaclaro,opacity=.45] (80,8) rectangle (106,-75);
\node[fp] (p2) at (93,0) {{\scriptsize P2}\\\textbf{Atenuar esforços}};
\node[fn] (e20) at (93,-19) {{\scriptsize E04}\\\textbf{Absorver impacto}};
\node[fn] (e21) at (93,-34) {{\scriptsize E05}\\\textbf{Isolar vibração}};
\node[fn] (e22) at (93,-49) {{\scriptsize E06}\\\textbf{Desviar compressão}};
\draw[energia] (93,-82)--(93,-75);
\node[porta,anchor=north,align=center,text width=25mm,font=\scriptsize] at (93,-84) {Choque, vibração e compressão};
\fill[cinzaclaro,opacity=.45] (107,8) rectangle (133,-75);
\node[fp] (p3) at (120,0) {{\scriptsize P3}\\\textbf{Conter fragmentos}};
\node[fn] (e30) at (120,-19) {{\scriptsize E07}\\\textbf{Envolver louça}};
\node[fn] (e31) at (120,-34) {{\scriptsize E08}\\\textbf{Reter fragmentos}};
\node[fn] (e32) at (120,-49) {{\scriptsize E09}\\\textbf{Vedar juntas}};
\draw[energia] (120,-82)--(120,-75);
\node[porta,anchor=north,align=center,text width=25mm,font=\scriptsize] at (120,-84) {Energia da quebra};
\fill[cinzaclaro,opacity=.45] (134,8) rectangle (160,-75);
\node[fp] (p4) at (147,0) {{\scriptsize P5}\\\textbf{Movimentar conjunto}};
\node[fn] (e40) at (147,-19) {{\scriptsize E13}\\\textbf{Receber força}};
\node[fn] (e41) at (147,-34) {{\scriptsize E14}\\\textbf{Transmitir força}};
\node[fn] (e42) at (147,-49) {{\scriptsize E15}\\\textbf{Estabilizar conjunto}};
\draw[energia] (147,-82)--(147,-75);
\node[porta,anchor=north,align=center,text width=25mm,font=\scriptsize] at (147,-84) {Trabalho do operador};
\fill[cinzaclaro,opacity=.45] (161,8) rectangle (187,-75);
\node[fp] (p5) at (174,0) {{\scriptsize P6}\\\textbf{Informar condição}};
\node[fn] (e50) at (174,-19) {{\scriptsize E16}\\\textbf{Expor superfície}};
\node[fn] (e51) at (174,-34) {{\scriptsize E17}\\\textbf{Indicar fechamento}};
\node[fn] (e52) at (174,-49) {{\scriptsize E18}\\\textbf{Exibir identidade}};
\fill[cinzaclaro,opacity=.45] (188,8) rectangle (214,-75);
\node[fp] (p6) at (201,0) {{\scriptsize P7}\\\textbf{Liberar louça}};
\node[fn] (e60) at (201,-19) {{\scriptsize E19}\\\textbf{Destravar apoios}};
\node[fn] (e61) at (201,-34) {{\scriptsize E20}\\\textbf{Abrir barreira}};
\node[fn] (e62) at (201,-49) {{\scriptsize E21}\\\textbf{Remover louça}};
\draw[energia] (201,-82)--(201,-75);
\node[porta,anchor=north,align=center,text width=25mm,font=\scriptsize] at (201,-84) {Trabalho do operador};
\draw[material] (p0.east) -- (p1.west);
\draw[material] (p1.east) -- (p2.west);
\draw[material] (p2.east) -- (p3.west);
\draw[material] (p3.east) -- (p4.west);
\draw[material] (p4.east) -- (p5.west);
\draw[material] (p5.east) -- (p6.west);
\draw[material] (0,0)--(p0.west);
\draw[material] (p6.east)--(244,0);
\draw[material] (214,-27)--(244,-27);
\node[porta,anchor=south west,text width=24mm] at (0,2) {Louça\\desembalada};
\node[porta,anchor=south west,text width=29mm] at (216,2) {Louça íntegra\\entregue};
\node[porta,anchor=south west,text width=29mm] at (216,-25) {Fragmentos retidos, se houver quebra};
\draw[energia] (93,8)--(93,17);
\node[porta,anchor=south,align=center,text width=48mm] at (93,19) {Energia dissipada};
\draw[sinal] (174,8)--(174,17);
\node[porta,anchor=south,align=center,text width=64mm] at (174,19) {Condição, fechamento e identidade};
\draw[material] (18,-104)--(32,-104);\node[legenda] at (35,-104) {Material};
\draw[energia] (88,-104)--(102,-104);\node[legenda] at (105,-104) {Energia};
\draw[sinal] (158,-104)--(172,-104);\node[legenda] at (175,-104) {Sinal};
\end{tikzpicture}
